\begin{table}[H]
\centering
\caption{\textbf{Behavioral Coherence vs Geometric Baselines (C1: Behavioral Coherence).} Comparison on the Michael Jordan circuit (same pinned nodes). Higher peak-token consistency and activation-pattern similarity indicate better functional grouping; lower sparsity consistency indicates activation concentration on fewer tokens. Silhouette and Davies-Bouldin are geometric quality indices (reported for context; baselines favor geometric compactness). \textsc{Pre-specified:} Behavioral metrics (peak-token consistency, activation-pattern similarity, sparsity consistency), equal node set across all methods. \textsc{Exploratory:} Cluster counts and thresholds were tuned to maximize respective quality indices for each method. \emph{Note:} Concept-aligned grouping achieves 2.3$\times$ higher token consistency and 5.8$\times$ higher activation similarity than cosine clustering, at the cost of lower geometric compactness (Silhouette 0.124 vs 0.707 for layer-adjacency). This trade-off reflects prioritization of functional coherence over spatial clustering. Replacement/Completeness scores are near-invariant across methods (Graph 0.69 $\rightarrow$ Subgraph 0.49; Completeness 0.90 $\rightarrow$ 0.81) because all three methods pin the same node set. These scores validate node coverage but do not discriminate grouping quality---behavioral metrics above provide the comparison.}
\label{tab:baselines}
\begin{tabular}{lrrr}
\toprule
\textbf{Metric} & \textbf{Concept-aligned} & \textbf{Cosine} & \textbf{Layer adjacency} \\
\midrule
Peak Token Consistency (↑)      & 0.425 & 0.183 & 0.301 \\
Activation Pattern Similarity (↑)& 0.762 & 0.130 & 0.415 \\
Sparsity Consistency (↓)         & 0.255 & 0.399 & 0.335 \\
Silhouette (↑)                   & 0.124 & -0.386 & 0.707 \\
Davies--Bouldin (↓)              & 1.30 & 1.58 & 0.49 \\
\bottomrule
\end{tabular}
\end{table}


